\chapter{Resultados}
\label{cap:resultados}

El capítulo de resultados presenta lo obtenido tras aplicar la metodología. No
debe limitarse a mostrar tablas o capturas: cada resultado debe ir acompañado de
una interpretación técnica y de una relación clara con los objetivos del TFG.

\section{Pruebas realizadas}
\label{sec:pruebas-realizadas}

Describe las pruebas, simulaciones, cálculos, mediciones o comparativas
realizadas. Indica las condiciones de partida, los parámetros utilizados, los
criterios de aceptación y cualquier incidencia relevante.

La Tabla~\ref{tab:resultados-ejemplo} muestra un formato simple para resumir
resultados cuantitativos. Ajusta las unidades, magnitudes y criterios al tipo de
trabajo realizado.

\begin{table}[!ht]
    \footnotesize
    \centering
    \caption[Ejemplo de resumen de resultados]{Ejemplo de resumen de resultados. Fuente: elaboración propia}
    \label{tab:resultados-ejemplo}
    \begin{tabular}{lccc}
        \toprule
        \textbf{Ensayo} & \textbf{Valor esperado} & \textbf{Valor obtenido} & \textbf{Error} \\
        \midrule
        Ensayo 1 & 10,00 & 9,80 & 2,0\% \\
        Ensayo 2 & 25,00 & 25,70 & 2,8\% \\
        Ensayo 3 & 50,00 & 49,10 & 1,8\% \\
        \bottomrule
    \end{tabular}
\end{table}

Cuando una tabla contenga muchos datos, resume en el cuerpo de la memoria los
valores más importantes y traslada el detalle completo a un anexo.

\section{Análisis de resultados}
\label{sec:analisis-resultados}

En esta sección se interpretan los resultados. Una buena discusión explica si
los valores obtenidos son coherentes, qué desviaciones aparecen, qué factores
las explican y qué consecuencias tienen para el diseño o la propuesta final.

Puedes estructurar el análisis por criterios:

\begin{itemize}
    \item Cumplimiento de requisitos técnicos.
    \item Comparación con el estado de la cuestión o con una solución de referencia.
    \item Robustez ante cambios en las condiciones de trabajo.
    \item Coste, mantenimiento, escalabilidad o facilidad de implantación.
\end{itemize}

\section{Discusión}
\label{sec:discusion}

La discusión conecta resultados y objetivos. Conviene responder explícitamente a
estas cuestiones:

\begin{enumerate}
    \item ¿Qué objetivo queda demostrado con cada resultado?
    \item ¿Qué resultados son más relevantes y por qué?
    \item ¿Qué hipótesis, requisitos o supuestos se confirman?
    \item ¿Qué limitaciones deben tenerse en cuenta al interpretar el trabajo?
\end{enumerate}

No presentes como concluyente un resultado que sólo se ha probado en condiciones
muy limitadas. Si una prueba es preliminar, dilo con claridad y propón cómo
debería ampliarse.

\section{Comparativa económica o técnica}
\label{sec:comparativa}

Si el TFG compara alternativas, este apartado puede incluir una valoración de
coste, rendimiento, dificultad de montaje, mantenimiento, consumo energético,
precisión, seguridad o cualquier otro criterio relevante.

\begin{table}[!ht]
    \footnotesize
    \centering
    \caption[Ejemplo de comparativa de alternativas]{Ejemplo de comparativa de alternativas. Fuente: elaboración propia}
    \label{tab:comparativa-alternativas}
    \begin{tabularx}{\textwidth}{lXXX}
        \toprule
        \textbf{Criterio} & \textbf{Alternativa A} & \textbf{Alternativa B} & \textbf{Alternativa C} \\
        \midrule
        Coste inicial & Bajo & Medio & Alto \\
        Complejidad & Baja & Media & Alta \\
        Escalabilidad & Media & Alta & Alta \\
        Riesgo técnico & Bajo & Medio & Medio \\
        \bottomrule
    \end{tabularx}
\end{table}

Justifica siempre las valoraciones cualitativas. Una tabla de comparación es
útil sólo si el texto explica de dónde proceden los criterios y cómo se han
ponderado.
